% TOP_PROBLEM: {"id": 30007089, "title": "Graph decomposition at the integral expectation threshold", "category_id": 2, "category": {"id": 2, "name": "combinatorics", "display_name": "Combinatorics", "description": "Counting problems, graph theory, discrete structures.", "slug": "combinatorics", "order_index": 2, "created_at": "2026-07-31T15:26:25.671Z"}, "status": "open"}
% FIRST_POSTED: 2026-10-06
% ATTEMPT_STATUS: solved
% OPENAI_CLAIM: {"schema_version":1,"source_repo":"https://github.com/openai/math","source_commit":"adc7f1241b42e322a6451854ab7e4b4c146bf78a","release_date":"2026-10-06","match":"full","resolution":"proved","independently_reviewed":false,"second_pass":"agreed","families":[{"family":"175","title":"Talagrand’s expectation thresholds, discrete convexity, and graph decompositions","subject":"Combinatorics","manuscripts":[{"title":"Graph Decompositions at the Integral Expectation Threshold","pdf_path":"preprints/Graph-Decompositions-at-the-Integral-Expectation-Threshold-October-5-2026/graph-threshold-decompositions.pdf","readme_path":"preprints/Graph-Decompositions-at-the-Integral-Expectation-Threshold-October-5-2026/README.md","date":"2026-10-05","theorem_ref":"preprints/Graph-Decompositions-at-the-Integral-Expectation-Threshold-October-5-2026/build/main.tex#thm:main"}],"lean":{"doc_path":"lean/docs/175.md","comparators":[],"declarations":[],"covers_main_theorem":false},"reasoning_trace":null}]}
\documentclass[11pt]{article}
\usepackage[margin=25mm]{geometry}
\usepackage{fontspec}
\setmainfont{Latin Modern Roman}
\usepackage{amsmath,amssymb,mathtools}
\usepackage{unicode-math}
\setmathfont{Latin Modern Math}
\usepackage{enumitem}
\usepackage{xurl,hyperref,bookmark}
\hypersetup{colorlinks=true,urlcolor=blue}
\setcounter{secnumdepth}{0}
\setlength{\parindent}{0pt}
\setlength{\parskip}{5pt}
\setlength{\emergencystretch}{3em}
\begin{document}
\section*{Graph decomposition at the integral expectation threshold}

\subsection{OpenAI's claim}
\textit{Graph pieces at the expectation threshold}: The paper claims a partition of any graph’s edge set into a universally bounded number of pieces whose containment thresholds are controlled by the graph’s integral expectation threshold. The pieces are fixed before host sampling, and their separate embeddings need not be compatible.

\subsection{Outline of the argument}
\textit{Graph pieces at the expectation threshold}: A covering input is converted into uniform neighborhood estimates for random hosts. Direct embedding lemmas handle the resulting structured graph pieces. The edge partition is assembled using those estimates so that a single universal number of pieces suffices, without requiring the piece embeddings to agree on their common vertices.

\subsection{Scope relative to this problem}
\textit{Graph pieces at the expectation threshold}: The manuscript’s result uses universal k>=2,L>=1 and all H on at most n vertices. Edge pieces may be empty/share vertices and their partition is fixed before the host is sampled. It compares each containment threshold with the integral q(H;n), not merely fractional q\_f.

With the hypotheses and conventions above, the claimed result addresses the whole mathematical assertion on this page.

\subsection{Formal verification}
Family 175 describes Lean coverage for the integral/fractional threshold comparison and the discrete-convexity statement; graph decomposition is a later separate consequence.

\textit{Graph pieces at the expectation threshold}: No manuscript-specific Comparator is listed in the pinned inventory.

The documented selection does not cover the entire manuscript assertion relevant here.

This site has not run the Lean development. These links describe OpenAI’s selected formal statements and should be read with their stated scope.

\subsection{Status caveat}
OpenAI reports that an internal model produced the manuscripts. Their arguments have not been independently verified as proofs by this site. Any portion outside the documented formal scope may contain mistakes, and the catalogue label reports the announced mathematical claim.

\subsection{Sources}
\begin{itemize}
\item OpenAI, \emph{Graph Decompositions at the Integral Expectation Threshold}, 2026-10-05; citation key \texttt{OAI:Graph-Decompositions-at-the-Integral-Expectation-Threshold-October-5-2026}. \href{https://github.com/openai/math/blob/adc7f1241b42e322a6451854ab7e4b4c146bf78a/preprints/Graph-Decompositions-at-the-Integral-Expectation-Threshold-October-5-2026/graph-threshold-decompositions.pdf}{PDF (GitHub)}; \href{https://github.com/openai/math/raw/adc7f1241b42e322a6451854ab7e4b4c146bf78a/preprints/Graph-Decompositions-at-the-Integral-Expectation-Threshold-October-5-2026/graph-threshold-decompositions.pdf}{PDF (direct)}; \href{https://github.com/openai/math/blob/adc7f1241b42e322a6451854ab7e4b4c146bf78a/preprints/Graph-Decompositions-at-the-Integral-Expectation-Threshold-October-5-2026/README.md}{Manuscript page and citation}; \href{https://github.com/openai/math/blob/main/preprints/Graph-Decompositions-at-the-Integral-Expectation-Threshold-October-5-2026/graph-threshold-decompositions.pdf}{Latest version}.
\item Family 175: \href{https://github.com/openai/math/blob/adc7f1241b42e322a6451854ab7e4b4c146bf78a/lean/docs/175.md}{Lean scope}.
\item \href{https://github.com/openai/math/blob/adc7f1241b42e322a6451854ab7e4b4c146bf78a/CONTENTS.md}{OpenAI catalogue entry}.
\end{itemize}
\end{document}
